\scanpage{387}\mspage{126}

\[
C_\infty(\underline t;a,b)\times\underline C_\infty(\underline t';b,c)\unsure{-|\longrightarrow}\underline C_\infty(\underline t'';a,c)
\]

donc une flèche dans $\mathrm{Hot}$

\[
\mathrm{comp}_{\tau_0}:C_\infty\times C_\infty- C_\infty
\]

Il faut voir qu'elle ne dépend pas du choix de $\tau_0$, mais si on a $\tau_0$ et $\tau'_0$, on les compare par

\diagram{\begin{tikzcd}[ampersand replacement=\&,row sep=0pt]
\& \tau_0\\
\tau_0\unsure{\times|*}\tau'_0\arrow[ur]\arrow[dr] \&\\
\&\tau'_0
\end{tikzcd}}

et il suffit donc de voir que s'il existe une flèche $\tau_0\longrightarrow\tau'_0$, alors $\mathrm{comp}_{\tau_0}=\mathrm{comp}_{\tau'_0}$. Mais une telle flèche induit \deleted{un diagramme commutatif} \inserted{des flèches} \deleted{\unsure{dans $\underline{\mathrm{Hom\,cof}}$}}, $\underline t_{\unsure{\mathrm I_0|\tau_0}}\longrightarrow\underline t_{\tau'_0}$, $\underline t'_{\tau_0}\longrightarrow\underline t'_{\tau'_0}$, $\underline t''_{\tau_0}\longrightarrow\underline t''_{\tau'_0}$, d'où un diagramme dans $\mathrm{Cat}$

\diagram{\begin{tikzcd}[ampersand replacement=\&,row sep=small,column sep=small,execute at end picture={\draw[thin] (leftupper.south west) -- (trialprime.north east); \draw[thin] (leftlower.south west) -- (rightupper.north east); \draw[thin] (leftmiddle.south west) -- (trial.north east);}]
\& |[alias=trial]|{\deleted{\unsure{t_{\tau_0}}}}\&|[alias=trialprime]|{\deleted{\unsure{t'_{\tau_0}}}}\&\\
\&|[alias=leftupper]|{\deleted{t''_{\tau_0}}}\arrow[rr,leftrightarrow]\arrow[dl]\arrow[dd]\&\&|[alias=rightupper]|{\deleted{t''_{\tau'_0}}}\arrow[dl]\arrow[dd,no head]\\
|[alias=leftmiddle]|{\deleted{\underline t_{\tau_0}}}\arrow[rr,leftrightarrow]\&\&{\deleted{\underline t_{\tau'_0}}}\&\\
\&|[alias=leftlower]|{\deleted{\underline t'_{\tau_0}}}\arrow[rr,leftrightarrow]\&\&{\deleted{t'_{\tau'_0}}}
\end{tikzcd}}

\diagram{\begin{tikzcd}[ampersand replacement=\&,column sep=small]
\deleted{\unsure{H}}C_\infty(\underline t_{\tau'_0};a,b)\times C_\infty(\underline t'_{\tau'_0};b,c)\arrow[r,"{\mathrm{comp}_{\tau'_0}}"]\arrow[d] \& C_\infty(t''_{\tau'_0};a,c)\arrow[d]\\
C_\infty(\underline t_{\tau_0};a,b)\times C_\infty(t'_{\tau_0};b,\unsure{c|a})\arrow[r,"{\mathrm{comp}_{\tau_0}}"] \& C_{\deleted{\unsure{\infty}}\infty}(\deleted{\unsure{t}}t''_{\tau_0},a,c)
\end{tikzcd}}

dont on vérifie tout de suite qu'il est commutatif. Or ce diagramme peut se réécrire dans \deleted{\unsure{Cat}} $\mathrm{Hot}$ comme

\diagram{\begin{tikzcd}[ampersand replacement=\&]
C_\infty(a,b)\times C_\infty(b,c)\arrow[r,"{\mathrm{comp}_{\tau'_0}}"]\arrow[d,"\mathrm{id}"]\&C_\infty(a,c)\arrow[d,"\mathrm{id}"]\\
C_\infty(a,b)\times C_\infty(b,c)\arrow[r,"{\mathrm{comp}_{\tau_0}}"]\&C_\infty(a,c),
\end{tikzcd}}

% archive sheet number: 386
